\documentclass[10pt,letterpaper,final]{report}
\usepackage[margin=.75in]{geometry}
\usepackage[utf8]{inputenc}
\usepackage{amsmath}
\usepackage{amsfonts}
\usepackage{amssymb}
\usepackage{amsthm}
\renewcommand\qedsymbol{$\blacksquare$}
\usepackage{enumerate}
\usepackage{enumitem}
\usepackage{hyperref}
\usepackage{pdfpages}
\usepackage{graphics}
\usepackage{graphicx}
\usepackage{tikz}
\usepackage{tikz-qtree}
\usepackage{forest}
\usepackage{float}
\usetikzlibrary{automata,arrows}
\usetikzlibrary{shapes.multipart, positioning}
\usepackage{mdframed}
\usepackage[
  backend=biber,
  style=numeric-comp,
  sorting=none
]{biblatex}
\addbibresource{references.bib}

\usepackage{minted}
\usemintedstyle{algol_nu}
\usepackage{xltabular}
\usepackage{pdflscape}
\usepackage{tcolorbox}
\usepackage{enumitem}
\usepackage{array}
\usepackage{todonotes}
\setuptodonotes{inline}  % render \todo notes as full-width blocks in the text flow, not margin boxes

\newcommand{\reason}[1]{&& \text{(#1)}}

% Based on template from COSC-417 at Towson University, originally authored by Marius Zimand

% To input a script, adapt the following:
% \inputminted[breaklines, breakanywhere, fontsize=\footnotesize, frame=single, fontsize=\footnotesize]{python}{script.py}

\begin{document}
\begin{center}
  \fbox{
    \vbox{
      \begin{flushleft}
        Jonathan Llovet \\  % authors' names
        EN.625.252.81 \\  %class
        2026-10-04 \\  % date
        Module 5 - Homework 5 \\ % ID
      \end{flushleft}
      \center{\Large{\textbf{Linear Algebra}}}
      %\end{mdframed}
    } % end vbox
  } % end fbox
  \vline
\end{center}

\medskip

\noindent\textbf{Directions}

Complete the following problems showing all your work. You may use a calculator to check your work, but should write out (or TeX up) all the steps of your solution. Unless otherwise specified, you may skip some steps of row reduction with a calculator, but state that you did so. Please upload your work as a single .pdf file in Canvas.

\medskip
\noindent\textbf{Problem 1:}
\medskip

Find a matrix $A$ whose inverse $A^{-1} = \begin{bmatrix}
    7 & 8 \\
    1 & 2
  \end{bmatrix}$

\medskip
\noindent\textbf{Answer:}
\medskip

\pagebreak

\noindent\textbf{Problem 2:}
\medskip

Suppose $A$ is invertible.

\begin{enumerate}
  \item Show or explain why $A^T A$ is also invertible for any matrix $A$.
  \item Show that $A^{-1} = (A^T A)^{-1} A^T$ by computing the product of $A$ and $(A^T A)^{-1} A^T$
\end{enumerate}

\medskip
\noindent\textbf{Answer:}
\medskip

\pagebreak

\noindent\textbf{Problem 3:}
\medskip

Let the matrices $A = \begin{bmatrix}
    3  & -2 & -2 \\
    -1 & 1  & 1  \\
    3  & -1 & -2
  \end{bmatrix}$ and $B = \begin{bmatrix}
    1  & a & 0  \\
    -1 & b & 1  \\
    2  & c & -1
  \end{bmatrix}$ be inverses of each other. Find the value of $b$.


\medskip
\noindent\textbf{Answer:}
\medskip

\pagebreak

\section*{Source Code}
The full source code, including LaTeX, tooling, other artifacts, and git history is available on my Github repository for the class here:
\begin{center}
  \url{https://github.com/jllovet/jhu-en-625-252-fa26-linear-algebra}
\end{center}

\printbibliography

\end{document}
